\documentclass{article}
\usepackage[T1]{fontenc}
\usepackage{iclr2027_conference}
\usepackage{times}
\usepackage{microtype}
\usepackage{hyperref}
\usepackage{url}
\usepackage{booktabs}
\usepackage{amsmath,amsfonts,amssymb}
\usepackage{array}
\usepackage{xcolor}
\usepackage{xspace}

\newcommand{\E}{\mathbb{E}}
\newcommand{\ind}{\mathbb{I}}
\newcommand{\cert}{\textsc{Certified}\xspace}
\newcommand{\bounded}{\textsc{Bounded}\xspace}
\newcommand{\notcert}{\textsc{Not-Certified}\xspace}
\newcommand{\tightlist}{\setlength{\itemsep}{2pt}\setlength{\parskip}{0pt}}

\title{When Is Planned Coordination Evaluable?\\
Estimand-Specific Receipts for Randomized\\
Multi-Agent Experiments}

\author{Anonymous Authors}

\begin{document}

\maketitle

\begin{abstract}
Randomized multi-agent experiments can bind assignment to a planned team while evaluating whatever execution record survives. A completion flag cannot establish that assignment, realized membership, target identity, terminal semantics, scorer inputs, protocol versions, and policy provenance jointly support the reported target. We study an estimand-indexed receipt gate that classifies three declared targets as \cert, \bounded, or \notcert from receipt-visible records and an external trust root. In the frozen synthetic Architecture stress census, the full gate produced 0/1,408,516 fatal false reports, retained 373,484/373,484 compatible cells, and had 0.9493686477 certified-cell coverage and 0.0499553870 null type-I error. A post-hoc descriptive extraction shows target-specific structural separation: on each relevant one-family fault row, the corresponding one-component ablation reported every unsupported case while the full gate reported none. This evidence is against a co-designed oracle; it does not establish independent empirical validity, real-building effects, live-model performance, or general diagnostic superiority. The contribution is a narrow agent-experiment binding specialization of existing support-identifiability principles and a reviewable synthetic consistency study.
\end{abstract}

\section{Problem}
\label{sec:problem}

Multi-agent evaluations usually describe an intervention at the level of a requested team or policy. The recorded endpoint, however, may have been generated by a different member set, joined to the wrong target, terminated under a stale semantic owner, scored with duplicated keys, interpreted under a different protocol version, or selected by a policy whose propensity was not retained. Cryptographic consistency within a mutated record does not answer the scientific question: which target does that record support?

The distinction matters because different causal targets require different evidence. Randomized assignment can support an intention-to-treat contrast despite execution noncompliance, while a controlled complete-bundle contrast requires the declared execution regime. A natural-policy value additionally requires policy provenance. A gate that accepts or rejects the run as a whole discards these distinctions; a permissive evaluator can instead report a precisely estimated but unsupported target.

We specialize existing support-identifiability principles to randomized multi-agent experiments. The specialization binds assignment, requested and realized execution, target identity, terminal semantics, scorer inputs, protocol versions, and policy provenance to each declared target. It is not a first receipt, trace system, observability framework, fault injector, causal estimator, generic identifiability theorem, or scientific-review agent. The finite observational-equivalence argument below is supporting mathematics for this declared model, not a general theorem claim.

Our study asks one question: under the declared randomized experiment and measurement model, which targets remain reportable after specific binding failures, and what descriptive inferential behavior accompanies those decisions in the frozen synthetic census? The answer is organized as a support lattice, a fail-closed gate, a synthetic stress benchmark with an independent transport boundary, and a bounded analysis of retained results.

\paragraph{Closest work and scope.}
Trace contracts and fault injection provide assurance machinery \citep{traceassurance2026}. AgentTelemetry supplies the OpenTelemetry-based public records used only for our transport check \citep{agenttelemetry2026}. TelemetrySuffBench separates failure detection from origin identification under insufficient telemetry \citep{telemetrysuffbench2026}. Most closely, Beyond Local Accuracy audits whether observation support separates policies with different evaluation targets \citep{beyondlocalaccuracy2026}. None of the trace or telemetry studies establishes our synthetic outcome model or comparative evaluator performance.

Our target definitions follow potential-outcome separation of intervention and outcome law \citep{rubin1974}; known-propensity weighting follows design-based sampling logic \citep{horvitzthompson1952}. Selective observation requires explicit missingness assumptions \citep{robinsrotnitzkyzhao1994}, partial identification and its uncertainty have an established literature \citep{manski1990,imbensmanski2004}, and endpoint misclassification is a classical measurement problem \citep{bross1954}. These foundations supply identification tools; they are not contributions of this paper. Our narrower object is an estimand-indexed binding contract for randomized multi-agent experiments and its frozen synthetic structural comparison.

\paragraph{Contributions.}
First, we declare three targets and an explicit support relation over receipt fields. Second, we implement a target-specific gate whose trust anchor remains outside the resealable artifact closure. Third, we report a post-hoc descriptive extraction of the pre-specified frozen evaluator configurations by target and one-family fault. Fourth, we preserve the public-census, V3-wrapper, external-transport, and post-hoc-custody boundaries rather than merging them into one validation claim.

\section{Declared Model and Estimands}
\label{sec:model}

\subsection{Units, assignment, execution, and measurement}

Let unit $u=(s,i)$ denote target $i$ in synthetic site $s$. The benchmark has eight sites and eight isolated targets per site. Within every site, four targets are assigned to each arm, so site clustering affects assignment and inference but the primary synthetic model assumes no interference across targets.

For each simulated replicate, $U$ is uniform over its fixed 64-position synthetic roster, and $\E$ in the three estimands is the finite-roster mean conditional on that replicate's generated potential outcomes and retained design objects. Monte Carlo performance summaries, unlike the estimands, additionally average over the declared site, outcome, balanced-assignment, and fault-selection streams across 2,000 replicates. Neither measure averages over real Architecture projects or deployed agents.

For arm $z\in\{0,1\}$, $A_u(z)$ is the requested action and complete specialist-member set, $E_u(z)$ is the realized action and member set, and
\begin{equation}
  C_u(z)=\ind\{E_u(z)=A_u(z)\}
\end{equation}
records complete compliance in every registered execution field. The external trust root commits to assignment $Z_u$, both requested regimes, the target join $B_u$, the current terminal owner and value $T_u^*$, the duplicate-free scorer map $S_u^*$, the protocol map $G_u^*$, and the policy record $P_u^*$. The policy record includes its identifier, action menu, information snapshot, assignment propensity, budget and cost rule, retry rule, and decision event.

Let $Y_u^*(z,e)\in[0,1]$ be the canonical endpoint under assignment $z$ and execution regime $e$. The measured endpoint $\widetilde Y_u$ is defined only when the target join, terminal semantics, scorer inputs, and protocol version establish the declared measurement link. This separation prevents a valid assignment from silently authenticating a different endpoint.

\subsection{Three declared targets}

The intention-to-treat target is
\begin{equation}
  \tau_{\mathrm{ITT}}
  = \E\!\left[Y_u^*\{1,E_u(1)\}-Y_u^*\{0,E_u(0)\}\right].
  \label{eq:itt}
\end{equation}
It follows randomized assignment and natural execution. Execution noncompliance therefore does not by itself remove ITT support, although broken assignment, target, terminal, scorer, or protocol bindings can do so.

The controlled complete-bundle target is
\begin{equation}
  \tau_{\mathrm{CB}}
  = \E\!\left[Y_u^*(1,e_1^C)-Y_u^*(0,e_0^C)\right],
  \label{eq:cb}
\end{equation}
where $e_z^C$ is the declared complete regime in arm $z$. A point report requires the execution receipt to establish that controlled regime or a separately specified partial-identification bound.

For a declared natural policy $\pi$, the policy value is
\begin{equation}
  \Psi_N(\pi)=\E\!\left[Y_u^*\{\pi(X_u),E_u(\pi(X_u))\}\right].
  \label{eq:policy}
\end{equation}
The retained estimator is site-averaged inverse-probability weighting with known propensity $1/2$. Policy identity, menu, snapshot, propensity, cost, retry, and decision-event provenance are part of its support.

\subsection{Support states}

For target $\theta\in\{\tau_{\mathrm{ITT}},\tau_{\mathrm{CB}},\Psi_N\}$ and receipt-visible record $R$, the gate returns:
\begin{description}
\tightlist
\item[\cert.] The frozen receipt conditions establish point support for $\theta$ under the declared model.
\item[\bounded.] The record does not establish a point target, but the declared assumptions and observed support establish a nontrivial valid interval.
\item[\notcert.] Neither a point target nor the declared bound is established by the receipt-visible record.
\end{description}
These labels are relative to the declared receipt-visible records absent additional validation data or an explicit measurement-error model. They do not assert that the physical or computational process is unknowable under every possible measurement regime.

\section{Estimand-Support Lattice and Gate}
\label{sec:gate}

Figure~\ref{fig:estimand-support-lattice} summarizes the main target-specific separations. Clean receipts support all three targets. A realized-member mismatch leaves assignment-based ITT and the natural-policy target supportable when their other bindings remain intact, but it removes point support for the controlled bundle. A policy-provenance mismatch is specific to $\Psi_N$. Measurement-link failures can affect every point target.

\begin{figure}[t]
\centering
\fbox{%
\begin{minipage}{0.96\linewidth}
\centering
\small
\setlength{\tabcolsep}{3pt}
\begin{tabular}{@{}>{\raggedright\arraybackslash}p{0.25\linewidth}
                  >{\centering\arraybackslash}p{0.20\linewidth}
                  >{\centering\arraybackslash}p{0.20\linewidth}
                  >{\centering\arraybackslash}p{0.20\linewidth}@{}}
\toprule
Receipt condition & $\tau_{\mathrm{ITT}}$ & $\tau_{\mathrm{CB}}$ & $\Psi_N(\pi)$ \\
\midrule
All bindings valid & \cert & \cert & \cert \\
Execution differs from request & \cert & \notcert & \cert \\
Target, terminal, scorer, or protocol link broken & \notcert & \notcert & \notcert \\
Design-specified partial bound only & \bounded & \bounded & \bounded \\
Policy provenance broken, other links valid & \cert & \cert & \notcert \\
Semantic invariance control & \cert & \cert & \cert \\
\bottomrule
\end{tabular}
\end{minipage}}
\caption{The estimand-support lattice. States are determined separately for each declared target from the external trust root and receipt-visible fields. A displayed \notcert means not certified under this declared record and model, not universal nonidentification.}
\label{fig:estimand-support-lattice}
\end{figure}

\subsection{Fail-closed decision procedure}

The evaluator receives anonymous same-format artifact bytes and the immutable trust root. It does not receive filenames, mutation labels, construction truth, or oracle output. Every evaluated artifact and internal reference may be altered and fully resealed; the attacker cannot alter the external root. The gate first parses a closed schema and canonical bytes, then checks root membership and the $Z,A,E,B,T,S,G,P$ bindings, and finally maps the surviving support to each estimand. A malformed field, unknown version, duplicate scorer key, stale terminal owner, or missing policy component fails closed for every target that requires it.

The oracle remains separate. It reads the clean construction and mutation ledgers only after evaluator decisions are frozen, determines which target is actually compatible with the mutated record, and scores false reportability and false nonreportability. Hash-only and schema-complete evaluators receive the same external root, so their comparison isolates semantics rather than access to stronger authentication.

\subsection{Component-relative nonidentification specialization}

Suppose a required binding component $c$ is absent from the projected observation $O^{-c}$. If two admissible laws $P,Q$ agree in distribution, $\mathcal L_P(O^{-c})=\mathcal L_Q(O^{-c})$, while $\theta(P)\ne\theta(Q)$ and $c$ is not recoverable from the retained fields or manifest, then no function of $O^{-c}$ point identifies $\theta$ over that declared world class. The eight finite witnesses cover one declared component--target example for each of $Z,A,E,B,T,S,G,P$; they do not cover every required component--target pair. Appendix~\ref{app:mathematical-details} states the assumptions, exact mapping, and proof. This is an application of observational equivalence, not a new generic theorem.

When the missing component is constrained by declared error rates or overlap structure, the same receipt can justify \bounded rather than \notcert. The gate emits a bound only when the assumptions and inputs named by the corresponding proposition are themselves bound to the root.

\section{Synthetic Stress Benchmark and Independent Boundary}
\label{sec:benchmark}

\subsection{Architecture stress setting}

The benchmark is a deterministic synthetic 8-by-8 roster used to stress assignment, membership, join, terminal, scorer, version, and policy bindings. ``Architecture'' is synthetic Architecture stress vocabulary; it does not denote a sampled population of buildings or a deployed design workflow. The benchmark is synthetic and provides no causal result about real buildings or evidence about a sampled Architecture population.

The public corpus contains 64 clean artifacts, 1,408 single-fault artifacts, 1,920 ordered two-fault compounds, and 192 semantic invariance controls, for 3,584 artifacts in total. Fault transforms operate upstream of evaluation and fully reseal the mutated closure. A separately authored 24-case pack contains harmful singles, harmful compounds, and invariance controls. Its truth was hidden from the evaluator and was not used to tune the public generator after reveal.

\subsection{Public artifact and external transport evidence}

The public census passed the retained gate contract. The initial V1 held-out preflight stopped before a scientific gate because two controls were no-ops. The correction replaced those no-op controls before V2. The original V2 held-out run was an invalid harness, not a case-level result; V3 was a vocabulary-only technical rerun. V2 expected lowercase decision tokens that were disjoint from the production uppercase vocabulary, so its aggregate mismatch contains no case-level semantic verdict. V3 changed only the wrapper vocabulary mapping and retained the pack and production pins. We do not describe V2 as rescued.

The external AgentTelemetry check uses 170 public records from the pinned revision \texttt{8ba0ae75} \citep{agenttelemetry2026}. The external AgentTelemetry check tests schema-identity transport only; it is not causal evidence and does not validate the synthetic Architecture setting. Five resealed binding tamper families were rejected after adaptation. No outcome from the source study was reinterpreted, and no language model or provider was called.

\begin{center}
\small
\textbf{External transport source pin}\\
commit \texttt{8ba0ae753d51cc2837f4ef2fa450e103c3be904a}
\end{center}

\subsection{Evaluators}

The frozen design contains four non-semantic configurations, six one-component ablations, and the full gate. Table~\ref{tab:evaluator-comparison} reports the decision-relevant one-family contrasts recovered from the frozen aggregate payload. This table is a post-hoc descriptive extraction from the pre-specified frozen configurations, not a confirmatory comparison backed by an external time-stamped commitment. Every row uses 24,000 replication-level decisions; a dash in the discussion means the omitted component is irrelevant to that target and the full gate remains reportable.

\begin{table*}[t]
\centering
\small
\setlength{\tabcolsep}{3pt}
\begin{tabular}{@{}c
                  >{\raggedright\arraybackslash}p{0.16\textwidth}
                  >{\raggedleft\arraybackslash}p{0.15\textwidth}
                  >{\raggedleft\arraybackslash}p{0.17\textwidth}
                  >{\raggedleft\arraybackslash}p{0.18\textwidth}
                  >{\raggedright\arraybackslash}p{0.22\textwidth}@{}}
\toprule
Fault & Target & Full FR & Full abstention & Ablated FR & Irrelevant-target behavior \\
\midrule
E & $\tau_{\mathrm{CB}}$ & 0/20{,}653 & 20{,}653/24{,}000 & 20{,}653/20{,}653 & ITT and $\Psi_N$: 0/24{,}000 abstention \\
B & All three & 0/20{,}636 & 20{,}636/24{,}000 & 20{,}636/20{,}636 & None \\
T & All three & 0/20{,}565 & 20{,}565/24{,}000 & 20{,}565/20{,}565 & None \\
S & All three & 0/20{,}619 & 20{,}619/24{,}000 & 20{,}619/20{,}619 & None \\
G & All three & 0/20{,}600 & 20{,}600/24{,}000 & 20{,}600/20{,}600 & None \\
P & $\Psi_N$ & 0/20{,}579 & 20{,}579/24{,}000 & 20{,}579/20{,}579 & ITT and $\tau_{\mathrm{CB}}$: 0/24{,}000 abstention \\
\bottomrule
\end{tabular}
\caption{Post-hoc target-specific fault-family comparison from frozen rows. FR is false reportability. ``Ablated'' removes the component named in the first column. The co-designed oracle supplies support labels, so this is structural stress evidence, not independent empirical validation.}
\label{tab:evaluator-comparison}
\end{table*}

\subsection{Frozen simulation and inference}

The retained design fixes 297 scenarios: the primary family plus design sensitivities, eight sites, balanced within-site assignment, and 2,000 replications per scenario. Separate configurations vary execution, binding, terminal, scoring, version, policy, and selection consequences. Random streams are Philox and separated by scenario and arm.

For each point-certified cell, the estimator is the mean of eight site-specific estimates. Confidence intervals use the standard error across sites and a two-sided Student $t$ critical value with seven degrees of freedom. The same interval-excludes-zero rule defines the retained null rejection rate; this is not an exact sharp-null randomization test. Bias, root mean squared error, coverage, type-I error, power, false sign, abstention, runtime, and bytes are defined only on their applicable denominators. Oracle-nonidentified targets receive \textsc{NA}, never zero, for point-error metrics.

Two authorized executions used the same fixed command apart from execution index and root. Their five scientific payloads were byte-identical. The second execution was a reproducibility check, not an analytic retry; no simulation was retuned or rescored after the results were revealed.

\section{Frozen Synthetic Results}
\label{sec:results}

The frozen evidence establishes a target-specific structural separation against the co-designed oracle in this synthetic fault census: the corresponding one-component ablation falsely reported every unsupported case in each relevant one-family row, whereas the full gate reported none. It does not establish independent empirical validity, general fault-diagnosis ability, causal recovery, or comparative inferential improvement.

Fatal-fault false reportability was $0/1{,}408{,}516=0$.

Compatible retention was $373{,}484/373{,}484=1$.

Certified-cell coverage was $354{,}574/373{,}484=0.9493686477$.

Null type-I error was 4,479/89,660, namely $0.0499553870$.

\begin{figure}[t]
\centering
\fbox{%
\begin{minipage}{0.94\linewidth}
\small
\textbf{Frozen full-gate point rates}\par\medskip
\setlength{\tabcolsep}{3pt}
\begin{tabular}{@{}>{\raggedright\arraybackslash}p{0.35\linewidth}
                  r
                  >{\raggedright\arraybackslash}p{0.36\linewidth}@{}}
Fatal false reportability & 0.00000 & \textcolor{black}{\rule{0.4pt}{1.25ex}} \\
Compatible retention & 1.00000 & \textcolor{black}{\rule{0.99\linewidth}{1.25ex}} \\
Certified-cell coverage & 0.94937 & \textcolor{black}{\rule{0.94\linewidth}{1.25ex}} \\
Null type-I error & 0.04996 & \textcolor{black}{\rule{0.05\linewidth}{1.25ex}} \\
\end{tabular}
\par\smallskip
\noindent\footnotesize Horizontal extent encodes the descriptive point rate on a common $[0,1]$ scale. Exact counts are reported in the text.
\end{minipage}}
\caption{Frozen descriptive point rates for the full estimand gate. Retention and coverage are conditional on compatible or certified cells; type-I error uses retained null cells. The shared simulation streams make the pooled indicators dependent.}
\label{fig:inference-consequence}
\end{figure}

The clean diagnostic has two denominators that answer different questions. Clean false rejection---failure to report a certified clean estimand-replication cell---was 0/54,000. Clean-null statistical rejection was separately 389/12,000, or 0.0324167. Conflating them would turn a reportability diagnostic into an inference claim, so we retain both labels and denominators.

The numerical cutoffs were pre-specified in the retained design narrative, but current custody does not independently timestamp them before result reveal. We therefore compare the point rates descriptively with the retained ranges---fatal false reportability at most 0.01, compatible retention at least 0.95, clean false rejection at most 0.01, certified-cell coverage from 0.925 to 0.975, and null type-I error from 0.03 to 0.07---without claiming externally verified preregistration.

Across applicable scenario rows, the coverage P10/median/P90 summaries were 0.930571/0.952500/1 for $\tau_{\mathrm{ITT}}$, 0.925520/0.955513/1 for $\tau_{\mathrm{CB}}$, and 0.931002/0.954000/1 for $\Psi_N$. Null-rejection P10/median/P90 was 0/0.0481477/0.0887162 for $\tau_{\mathrm{ITT}}$ and 0/0.0455000/0.100000 for $\tau_{\mathrm{CB}}$. These empirical scenario-row quantiles are heterogeneity summaries, not confidence intervals. No defensible pooled Monte Carlo standard error is recoverable because replicate-level joint sufficient statistics and cross-estimand and cross-configuration covariance were not persisted.

The public artifact census had zero harmful false reportability and zero clean or invariance false nonreportability. The disclosed V3 vocabulary-only technical rerun retained the harmful-fault and invariance behavior of its 24-case pack; the invalid V2 harness supplied no case-level result. Exact paired-world law equalities, target inequalities, declared bounds, terminal-misclassification equalities, and selection counterexamples were checked at the frozen tolerance. These are implementation-consistency facts against co-designed witnesses, not independent validation. The strongest independent held-out/external macro comparison was not preserved in the frozen payload, so we make no claim that its pre-specified margin criterion was met.

\section{Limitations and Validity Boundaries}
\label{sec:limitations}

\paragraph{Synthetic domain.}
The roster, outcomes, faults, and scenario census are synthetic. There is no causal real-building result, no representative sample of Architecture practice, and no evidence that the numerical rates transport to deployed workflows. The Architecture terms are synthetic Architecture stress vocabulary only.

\paragraph{No live model or provider result.}
No live large language model, model provider, paid API, human-derived hidden case, or production agent service was evaluated. The paper therefore makes no claim about a model's reasoning quality, a provider's reliability, or the operational benefit of a multi-agent system.

\paragraph{External transport is not external causal validity.}
The independently authored AgentTelemetry records show that the schema-identity adapter and resealed binding checks operate on a public trace shape. They do not validate our synthetic outcome model, identify an effect in the AgentTelemetry study, or establish causal transport to another domain.

\paragraph{Declared records and assumptions.}
\notcert\ and nonidentified refer only to the declared receipt-visible records and world class. Extra validation measurements, replicated outcomes, or an explicit measurement-error model could change support. The finite witnesses do not establish a universal impossibility theorem.

\paragraph{Post-hoc custody.}
The registered roots were made read-only and independently hashed only after scoring. Two matching roots are retained, but the post-hoc custody correction cannot cryptographically prove the historical absence of a deleted third execution. Current hashes also cannot prove that no prior mutation occurred before the freeze. We therefore claim byte agreement of the two retained roots, not a cryptographic proof of the entire execution history.

\paragraph{Protocol deviation.}
The invalid V2 held-out wrapper and the V3 vocabulary-only technical rerun weaken the cleanliness of the held-out narrative even though V2 produced no case-level semantic result and production bytes remained pinned. Readers should treat the public census, retained simulation, and external transport as distinct evidence streams rather than merging them into a single validation score.

\section{Conclusion}
\label{sec:conclusion}

A completion marker is not a scientific support certificate. For a declared randomized multi-agent experiment, assignment-based, controlled-execution, and natural-policy targets depend on different combinations of execution and provenance evidence. An estimand-indexed receipt makes those requirements explicit and permits compatible targets to remain reportable while unsupported point claims fail closed.

In the frozen synthetic stress study, the full gate had zero observed fatal false reports and the corresponding one-component ablation falsely reported every unsupported case in each relevant one-family row. That target-specific structural contrast is limited to the co-designed synthetic census and receipt-visible model. The external public trace check establishes schema transport only, and the study contains no causal result about real Architecture work, live language models, or providers. The defensible novelty is the agent-experiment binding specialization and its auditable synthetic behavior, not a generic support theorem, general diagnostic system, or comparative inferential advance.

\section*{Reproducibility and Custody Statement}

The retained design narrative fixes the scenario order, evaluator profile, numerical cutoffs, estimators, confidence-interval rule, and random-stream identities, but the current chain does not independently prove an external pre-result timestamp. The scientific payloads of the two retained roots are byte-identical. Results in this manuscript were copied from the frozen retained-results report; no simulation rerun, rescore, post-reveal tuning, or protected-data access was used for manuscript preparation. The custody limitation above remains part of the evidence record.

\section*{ICLR 2027 AI Use Statement}

AI assistance was used during conceptual and mathematical development, study design, implementation, interpretation, literature work, and manuscript drafting and editing. Humans reviewed the cited primary sources, mathematical statements, implementation behavior, frozen outputs, disclosures, and final wording. Humans remain responsible for every claim, source, proof, implementation decision, result interpretation, and the final text. No AI-generated text or analysis is treated as an independent correctness certificate, and no unverified model or provider identity is asserted here.

\bibliography{references}
\bibliographystyle{iclr2027_conference}

\appendix

\section{Mathematical Details and Assumptions}
\label{app:mathematical-details}

\subsection{Assumptions and support lattice}

The displayed lattice is relative to the following declared model. The external manifest correctly binds the 64-position synthetic roster and named design objects; outcomes lie in $[0,1]$; targets do not interfere within the synthetic construction; assignment is balanced within site; and consistency links each declared assignment and execution regime to its corresponding potential outcome. A measured endpoint is usable only when the $B,T,S,G$ measurement link is valid. The controlled-bundle target additionally requires $A,E$, and the natural-policy target additionally requires the complete $P$ record, randomized conditional exchangeability, and positivity. Bounds require their named error-rate or overlap inputs; absent those inputs the gate does not substitute a point value.

For a receipt $R$ and target $\theta$, let $\mathcal S_\theta(R)$ be the set of admissible target values under these assumptions and the root-bound fields. The lattice labels mean $|\mathcal S_\theta(R)|=1$ for \cert, a nontrivial proper interval for \bounded, and no certified point or declared interval for \notcert. These are reportability states under the declared observation model, not assertions about every possible auxiliary measurement.

\subsection{Component-relative nonidentification}

\paragraph{Proposition 1.}
For target $\theta$, component $c$, model class $\mathcal M_\theta$, and projected observation $O^{-c}$, suppose $c$ is not almost surely recoverable from $O^{-c}$ or the authenticated manifest. If $P,Q\in\mathcal M_\theta$ satisfy $\mathcal L_P(O^{-c})=\mathcal L_Q(O^{-c})$ but $\theta(P)\ne\theta(Q)$, then $\theta$ is not point identified from $\sigma(O^{-c})$ over $\mathcal M_\theta$.

\paragraph{Proof.}
If a measurable functional $h(O^{-c})$ identified $\theta$, equality in law would give the same distribution, and hence the same constant identified value, under $P$ and $Q$. This contradicts $\theta(P)\ne\theta(Q)$.

The exact finite witness roster is $Z\mapsto\tau_{\mathrm{ITT}}$, $A\mapsto\tau_{\mathrm{CB}}$, $E\mapsto\tau_{\mathrm{CB}}$, $B,T,S,G\mapsto\tau_{\mathrm{ITT}}$, and $P\mapsto\Psi_N$. Each pair has equal projected law, unequal target value, and a manifest that does not recover the omitted component. This roster proves one selected example per component, not every component--target pair.

\subsection{Overlap-aware bounded contrast error}

\paragraph{Proposition 2.}
For arm $z$, let nonnegative analysis weights $\alpha_{zu}$ sum to one. Let $M_{zu}\in\{0,1\}$ mark target replacement, $q_{zu}\in[0,1]$ be omitted intended positive key mass, and $d_{zu}=\operatorname{TV}(w^R_{zu},\widetilde w_{zu})\in[0,1]$ be residual total variation between normalized retained-key and realized weights. With intended and measured outcomes in $[0,1]$, fixed arm denominators, no unit deletion, no cross-arm movement, and no negative weights, define
\begin{equation}
b_z=\sum_u\alpha_{zu}\left[M_{zu}+(1-M_{zu})\min\{1,q_{zu}+d_{zu}\}\right].
\end{equation}
Then $|\widetilde\tau-\tau|\le b_1+b_0\le2$.

\paragraph{Proof.}
Target replacement is charged first and contributes at most one because both outcomes lie in $[0,1]$. Conditional on no replacement, omitting positive key mass changes the normalized score by at most $q_{zu}$, and residual reweighting changes the expectation of a $[0,1]$ score by at most its total-variation distance $d_{zu}$. The triangle inequality and the unit-range cap give the bracketed term without double counting overlap. Weighting within arm and applying the triangle inequality to the treated-minus-control contrast yields $b_1+b_0$; each arm bound is at most one. Denominator changes, deletion, cross-arm movement, negative weights, or scores outside $[0,1]$ are outside this proposition and do not receive this bound.

\subsection{Terminal misclassification and selection}

\paragraph{Proposition 3.}
For a binary endpoint with arm mean $\mu_z$, sensitivity $se_z$, specificity $sp_z$, $\kappa_z=se_z+sp_z-1$, no other binding or selection error, and identical standardized covariate composition,
\begin{equation}
\widetilde\tau=sp_0-sp_1+\kappa_1\mu_1-\kappa_0\mu_0.
\end{equation}
With common nondifferential $se,sp$, this reduces to $\widetilde\tau=(se+sp-1)\tau$.

\paragraph{Proof.}
In arm $z$, $\widetilde\mu_z=se_z\mu_z+(1-sp_z)(1-\mu_z)=(1-sp_z)+\kappa_z\mu_z$. Subtracting arm zero from arm one gives the displayed expression. Common rates cancel the additive terms and factor the true contrast. For example, $\mu_0=0.4$, $\mu_1=0.6$, $(se_0,sp_0)=(1,1)$, and $(se_1,sp_1)=(0.5,1)$ change the true contrast $0.2$ to $-0.1$.

\paragraph{Selection boundary.}
If the canonical outcome remains observed when a receipt fails, the randomized target uses the full sample. Under receipt missingness at random given observed $X,Z$ and positivity, a declared standardization or inverse-probability functional can identify the target \citep{robinsrotnitzkyzhao1994}. Outcome-dependent missingness after conditioning generally leaves the target only partially identified without further restrictions \citep{manski1990,imbensmanski2004}. Conditioning on post-treatment compliance does not by itself identify a per-protocol causal effect. These MCAR, MAR, and MNAR statements are assumptions and identification boundaries, not properties certified merely by receipt parsing.

\end{document}
